\documentclass[11pt,a4paper]{article}
\usepackage[T1]{fontenc}
\usepackage[utf8]{inputenc}
\usepackage[margin=1in]{geometry}
\usepackage{lmodern}
\usepackage{amsmath,amssymb,graphicx,booktabs,tabularx,array}
\usepackage{microtype}
\usepackage{caption}
\usepackage{placeins}
\usepackage[hidelinks]{hyperref}
\usepackage{float}

\captionsetup{font=small,labelfont=bf}
\setlength{\emergencystretch}{2em}
\newcolumntype{Y}{>{\raggedright\arraybackslash}X}

\newcommand{\figpath}{../figures/Whakaari}

\newcommand{\pending}[1]{\par\smallskip\noindent\textit{To complete: #1}\par\smallskip}

\title{Ordinal Pattern Analysis of RSAM Time Series at Whakaari, 2014--2019}
\author{Rasika Dilhani}
\date{\today}

\begin{document}
	\maketitle
	\begin{abstract}
		This study investigates ordinal pattern features for retrospective detection of unusual seismic-amplitude behavior at Whakaari, New Zealand. RSAM records from WIZ and WSRZ cover the intended period July 2014-December 2019, giving corrected modeling datasets of 335 and 333 weekly windows, respectively, including four eruption-labeled windows at each station. Three nested feature sets (thirteen, seven, and four descriptors) are compared across three analysis methods: K-means partitioning, DBSCAN clustering, and a class-weighted radial-basis SVM evaluated under four class-imbalance treatments (no balancing, oversampling, undersampling, SMOTE). Neither K-means nor DBSCAN recovers a partition resembling the three Volcanic Alert Level (VAL) classes (Adjusted Rand Index below 0.06 in every configuration at both stations). Baseline SVM classification (seven features, no balancing) reaches an overall accuracy of 0.627 at both stations, but a one-versus-rest balanced accuracy of only 0.365 at WIZ and 0.530 at WSRZ; WSRZ recovers two of four eruption windows under most feature-set and balancing combinations, while WIZ recovers at most one of four, and only at a large cost to background-class accuracy. These exploratory results characterize unusual feature-space behavior and quantify the limits of the tested clustering and classification configurations, but do not establish reliable eruption recognition or forecasting.
	\end{abstract}
	
\section{Introduction}
Real-time Seismic Amplitude Measurement (RSAM) summarizes seismic amplitude over specified intervals and supports volcano monitoring~\cite{Endo1991}. Amplitude summaries provide a useful overview, while changes in the temporal ordering of observations may supply complementary information. Ordinal-pattern analysis represents this ordering through distributions of ranked observation sequences~\cite{Bandt2002}.

This study examines whether ordinal features identify unusual weekly RSAM windows at Whakaari and how these windows relate to Volcanic Alert Level (VAL) records. We analyze WIZ and WSRZ separately and compare them over a common study period. Two unsupervised methods, K-means partitioning and DBSCAN, are used to ask whether the ordinal-pattern feature space separates into groups resembling the three VAL classes; a class-weighted support vector machine (SVM) addresses the corresponding supervised classification problem. All three methods are additionally compared across nested feature sets of thirteen, seven, and four descriptors, and the SVM is compared across four treatments for the pronounced class imbalance between background, moderate-unrest, and eruption weeks. The objective is retrospective characterization and detection, with forecasting claims requiring a separate evaluation using only information available before each event.

\section{Study Area and Data Overview}\label{sec:data}
\subsection{Records and Scope}
The study concerns Whakaari, New Zealand, and RSAM records identified by stations WIZ and WSRZ. The intended analysis interval is 1 July 2014-31 December 2019. Figure axes are labeled in New Zealand local time. 	

\subsection{Long-Term RSAM Records}
Figures~\ref{fig:full_wiz} and~\ref{fig:full_wsrz} introduce the full records before the feature analysis. Both show a relatively low baseline punctuated by transient peaks, with differences in amplitude and timing between stations. 

\begin{figure}[H]
	\centering
	\includegraphics[width=\linewidth]{\figpath/WIZ_2014to2019_highlighted.pdf}
	\caption{WIZ RSAM record for the 2014-2019 interval.}
	\label{fig:full_wiz}
\end{figure}	

\begin{figure}[H]
	\centering
	\includegraphics[width=\linewidth]{\figpath/WSRZ_2014to2019_highlighted.pdf}
	\caption{WSRZ RSAM record for the 2014-2019 interval.}
	\label{fig:full_wsrz}
\end{figure}	


\subsection{Selected Short-Interval Views}
Figures~\ref{fig:apr2016}-\ref{fig:dec2019interval} compare the two stations around the dates selected in the supplied plots. Green represents WIZ and purple represents WSRZ. 

\begin{figure}[H]
	\centering
	\includegraphics[width=\linewidth]{\figpath/zoom_28_Apr_2016_eruption.pdf}
	\caption{RSAM comparison around the 28 April 2016 annotation.}
	\label{fig:apr2016}
\end{figure}

\begin{figure}[H]
	\centering
	\includegraphics[width=\linewidth]{\figpath/zoom_13_15_Sep_2016_eruption.pdf}
	\caption{RSAM comparison around the 13-15 September 2016 annotation. The isolated WIZ peak precedes the highlighted interval.}
	\label{fig:sep2016}
\end{figure}

\FloatBarrier
\begin{figure}[H]
	\centering
	\includegraphics[width=\linewidth]{\figpath/zoom_9_Dec_2019_explosive_eruption.pdf}
	\caption{RSAM comparison around 9 December 2019.}
	\label{fig:dec2019onset}
\end{figure}

\begin{figure}[H]
	\centering
	\includegraphics[width=\linewidth]{\figpath/zoom_9_12_Dec_2019_continued_eruptive_activity.pdf}
	\caption{RSAM comparison with the 9-12 December interval.}
	\label{fig:dec2019interval}
\end{figure}	


\section{Literature Review}

\subsection{Information-Theoretic Measures in Volcanic Seismology}

Information-theoretic measures quantitatively describe the variability and statistical structure of seismic signals. In volcanic monitoring, researchers have investigated these measures as complementary tools for identifying changes in seismic activity associated with eruptive processes. Telesca et al.~\cite{Telesca2014} applied the Fisher-Shannon method to continuous seismic recordings collected before and during the 2011-2012 El Hierro eruption in the Canary Islands. Their work illustrates how complementary information measures can characterize the evolution of volcanic seismic signals.

Rey-Devesa et al.~\cite{ReyDevesa2023} investigated Shannon entropy in continuous seismic recordings from Volc\'an de Colima, Mexico. They reported decreases in entropy preceding two major eruptive episodes and examined entropy minima as indicators of smaller explosions, supported by visual monitoring observations. These findings suggest that entropy can provide information relevant to tracking volcanic activity. However, entropy variations alone do not establish a unique physical mechanism or demonstrate deterministic conduit behavior. Furthermore, Shannon entropy derived from seismic waveform statistics should be distinguished from permutation entropy, which describes the probability distribution of temporal ordinal patterns.

\subsection{Ordinal Patterns and Permutation Entropy}

Bandt and Pompe~\cite{Bandt2002} introduced permutation entropy as a method for characterizing time series through the relative ordering of observations. For a selected embedding dimension $D$ and time delay $\tau$, sequences of observations are represented by ordinal patterns. The relative frequencies of these patterns define a probability distribution, from which permutation entropy is calculated. This representation captures aspects of temporal organization without requiring a specified parametric distribution for the signal amplitudes. Although the method is computationally efficient and can be useful for noisy observations, its performance depends on the embedding parameters, the number of available observations, and how tied values are treated. Its application to non-stationary signals also requires analysis windows appropriate to the timescales under investigation.

Statistical complexity complements entropy by capturing a distribution's departure from a reference state. Lamberti et al.~\cite{lamberti2004intensive} developed an intensive statistical complexity measure using Jensen--Shannon divergence to quantify disequilibrium. In this framework, normalized Shannon entropy $H$ and normalized disequilibrium $Q_J$ are combined as $C=H Q_J$. When calculated from ordinal-pattern distributions, these measures provide a joint representation of temporal pattern diversity and departure from a uniform distribution in the $H \times C$ plane.

\subsection{State Transitions in Geysers and Implications for Volcanic Monitoring}

Sudibyo et al.~\cite{Sudibyo2022} investigated permutation entropy in seismic recordings from Strokkur geyser, Iceland. Their analysis of 63 eruptive cycles identified recurring entropy variations associated with different eruption-cycle phases. They reported a systematic increase in permutation entropy during the final approximately 15 s before eruption. They found that permutation entropy outperformed seismic signal strength or quiescence for forecasting eruptions in the studied setting. Their work also examined how parameter selection and station distance affected the resulting entropy measurements.

These findings provide empirical motivation for investigating ordinal pattern measures as indicators of changes in volcanic and hydrothermal activity. However, transferring results from Strokkur to Whakaari requires separate evaluation. Repeated geyser eruptions observed through high-frequency seismic recordings differ from volcanic activity examined through aggregated RSAM time series over weekly windows. Consequently, the direction, timing, and physical interpretation of entropy changes cannot be assumed to remain consistent across these systems.

In the present study, the Strokkur investigation is relevant because it shows that temporal ordinal structure can contain information about changes in an eruptive system. This supports examining ordinal feature trajectories in Whakaari records while distinguishing retrospective anomaly detection from eruption forecasting. A window containing eruptive or post-eruptive observations may help characterize an eruption-related change, but it cannot establish a pre-eruptive signal.

\subsection{Motivation for Multi-Station Feature Analysis}
The reviewed studies motivate examining information measures in long-term RSAM records, but do not directly establish the performance of the present combination of ordinal features at Whakaari. Entropy, complexity, disequilibrium, and variance descriptors are overlapped. Feature selection is therefore used to seek a compact representation for clustering and classification. Comparing WIZ and WSRZ can reveal whether detected changes are shared across stations, while retaining station-specific differences in data quality and signal characteristics.

\section{Methodology} 

\subsection{Preprocessing and Weekly Windows} 

The feature-generation description specifies 10-minute RSAM observations and seven-day windows. The supplied clustering script reads precomputed weekly features and does not construct windows; therefore, this script does not verify the step size. A complete window on a regular UTC grid contains $N=7\times144=1008$ observations. The window duration and step are separate settings: a seven-day duration does not imply a seven-day step. 

For the full intended interval, a daily step would produce about 2000 candidate windows before exclusions, whereas the corrected modeling output contains 335 WIZ and 333 WSRZ windows. The difference depends on the feature-generation settings and exclusion counts, not an assumed daily step. Seven local calendar days can also differ from 168 hours at daylight-saving transitions; the time convention must be explicit. 


\subsection{Ordinal Pattern Distribution}

For a regular time series $X=\{x_t\}_{t=1}^{N}$, ordinal patterns are constructed from the embedding vectors
\begin{equation}
	\mathbf{x}_t =
	(x_t,x_{t+\tau},\ldots,x_{t+(D-1)\tau}),
\end{equation}
where $D$ denotes the embedding dimension and $\tau$ is the embedding delay. Each vector is mapped to the permutation that describes the relative ordering of its components.

In this study, the embedding parameters are $D=5$ and $\tau=1$ observation. Therefore, the number of possible ordinal patterns is
\begin{equation}
	k=D!=5!=120.
\end{equation}
For a complete seven-day window containing $N=1008$ observations, the number of available embedding vectors is
\begin{equation}
	M=N-(D-1)\tau=1004.
\end{equation}

Let $n_\ell$ denote the number of occurrences of ordinal pattern $\ell$. The empirical ordinal-pattern probability distribution is then
\begin{equation}
	\widehat{p}_\ell=\frac{n_\ell}{M},
	\qquad
	\ell=1,\ldots,k,
\end{equation}
and
\begin{equation}
	\widehat{\mathbf{P}}
	=
	\{\widehat{p}_\ell\}_{\ell=1}^{k},
	\qquad
	\sum_{\ell=1}^{k}\widehat{p}_\ell=1.
\end{equation}

If embedding vectors containing invalid or missing observations are excluded, $M$ denotes the number of valid vectors actually used to estimate the distribution. A consistent tie-breaking rule is applied when equal values occur within an embedding vector. Because the discrete Fisher Information Measure depends on differences between adjacent pattern probabilities, a fixed ordering of the ordinal patterns is also retained throughout the analysis.

\subsection{Information-Theoretic Quantifiers}

For each empirical ordinal-pattern probability distribution
$\widehat{\mathbf{P}}=\{\widehat{p}_\ell\}_{\ell=1}^{k}$,
a collection of global and local information-theoretic measures is calculated. These measures characterize different properties of the temporal dynamics represented by the ordinal-pattern distribution. Entropy-based measures describe global uncertainty and disorder, whereas the Fisher Information Measure describes local variation among neighbouring probabilities.

\subsubsection{Permutation Entropy}

Shannon entropy provides a measure of the uncertainty contained in the ordinal-pattern probability distribution. The Shannon entropy of
$\widehat{\mathbf{P}}$ is
\begin{equation}
	S(\widehat{\mathbf{P}})
	=
	-\sum_{\ell=1}^{k}
	\widehat{p}_\ell\log \widehat{p}_\ell,
\end{equation}
where terms with $\widehat{p}_\ell=0$ are defined to contribute zero.

The maximum Shannon entropy occurs for the uniform distribution and is equal to $\log k$. The normalized Shannon entropy, commonly referred to as permutation entropy in the ordinal-pattern framework, is therefore defined as
\begin{equation}
	H_{S}(\widehat{\mathbf{P}})
	=
	-\frac{1}{\log k}
	\sum_{\ell=1}^{k}
	\widehat{p}_\ell\log \widehat{p}_\ell.
\end{equation}
Consequently,
\begin{equation}
	0\leq H_{S}(\widehat{\mathbf{P}})\leq1.
\end{equation}

Low permutation entropy indicates that only a small number of ordinal patterns dominate the distribution, corresponding to a comparatively regular or ordered signal. Values close to one indicate a distribution approaching uniformity and therefore greater disorder or unpredictability. Entropy alone, however, cannot distinguish complete randomness from more structured forms of dynamical organization. For this reason, additional information and complexity measures are considered.


\subsubsection{Tsallis Entropy}

For $q\in\mathbb{R}\setminus\{1\}$, the Tsallis entropy is defined as
\begin{equation}
	S_{T}^{(q)}(\widehat{\mathbf{P}})
	=
	\frac{1}{q-1}
	\left(
	1-\sum_{\ell=1}^{k}
	\widehat{p}_\ell^{\,q}
	\right).
\end{equation}
The Shannon entropy is recovered in the limit $q\rightarrow1$.

The maximum Tsallis entropy occurs for the uniform distribution and is
\begin{equation}
	S_{T,\max}^{(q)}
	=
	\frac{1-k^{1-q}}{q-1}.
\end{equation}
Hence, the normalized Tsallis entropy used in the analysis is
\begin{equation}
	H_{T}^{(q)}(\widehat{\mathbf{P}})
	=
	\frac{
		S_{T}^{(q)}(\widehat{\mathbf{P}})
	}{
		S_{T,\max}^{(q)}
	}.
\end{equation}

The parameter $q$ controls the relative contribution of common and rare ordinal patterns and therefore provides an alternative characterization of the ordinal-pattern distribution to Shannon entropy.

\subsubsection{R\'enyi Entropy}

For $q>0$, $q\neq1$, the R\'enyi entropy is defined as
\begin{equation}
	S_{R}^{(q)}(\widehat{\mathbf{P}})
	=
	\frac{1}{1-q}
	\log
	\left(
	\sum_{\ell=1}^{k}
	\widehat{p}_\ell^{\,q}
	\right).
\end{equation}
As with Tsallis entropy, Shannon entropy is obtained in the limit $q\rightarrow1$.

The maximum R\'enyi entropy is $\log k$, obtained for the uniform distribution. The normalized R\'enyi entropy is therefore
\begin{equation}
	H_{R}^{(q)}(\widehat{\mathbf{P}})
	=
	\frac{
		S_{R}^{(q)}(\widehat{\mathbf{P}})
	}{
		\log k
	}.
\end{equation}

Different values of $q$ change the relative emphasis placed on ordinal patterns with high and low probabilities and provide complementary descriptions of the underlying temporal dynamics.


\subsubsection{Fisher Information Measure}

Entropy measures characterize global properties of the ordinal-pattern distribution. To complement these measures with a descriptor that is sensitive to local changes in the probability distribution, we adopt the discrete Fisher Information Measure proposed by Rey et al.~\cite{Rey2024Asymptotic}:
\begin{equation}
	F(\widehat{\mathbf{P}})
	=
	4
	\sum_{\ell=1}^{k-1}
	\left(
	\sqrt{\widehat{p}_{\ell+1}}
	-
	\sqrt{\widehat{p}_{\ell}}
	\right)^2.
\end{equation}

Unlike entropy, which depends on the overall distribution of probabilities, the Fisher Information Measure depends on differences between neighbouring probabilities under the adopted ordinal-pattern ordering. It therefore captures local gradients and fine-scale changes in the probability distribution that may not be evident from entropy measures alone.


\subsubsection{Disequilibrium}
\label{disequilibrium}

The entropy and Fisher Information measures characterize complementary properties of the ordinal-pattern distribution. However, an additional quantity is required to measure the departure of the observed distribution from complete randomness. In the ordinal-pattern framework, the natural equilibrium reference is the uniform distribution
\begin{equation}
	\mathbf{U}
	=
	\left(
	\frac{1}{k},
	\frac{1}{k},
	\ldots,
	\frac{1}{k}
	\right).
\end{equation}

Following the disequilibrium concept introduced by Lopez-Ruiz et al.~(\cite{lopez1995statistical}), the departure from the uniform distribution is quantified using the Jensen--Shannon divergence. Let
\begin{equation}
	S(\mathbf{P})
	=
	-\sum_{\ell=1}^{k}p_\ell\log p_\ell
\end{equation}
denote the unnormalized Shannon entropy. The Jensen--Shannon divergence between the empirical distribution $\widehat{\mathbf{P}}$ and the uniform distribution $\mathbf{U}$ is
\begin{equation}
	J(\widehat{\mathbf{P}},\mathbf{U})
	=
	S\left(
	\frac{\widehat{\mathbf{P}}+\mathbf{U}}{2}
	\right)
	-
	\frac{1}{2}S(\widehat{\mathbf{P}})
	-
	\frac{1}{2}S(\mathbf{U}),
\end{equation}
where
\begin{equation}
	S(\mathbf{U})=\log k.
\end{equation}

The maximum possible Jensen--Shannon divergence from the uniform distribution occurs when all probability mass is concentrated on a single ordinal pattern and is given by
\begin{equation}
	J_{\max}
	=
	\log 2
	+
	\frac{1}{2}\log k
	-
	\frac{k+1}{2k}\log(k+1).
\end{equation}

The normalized disequilibrium is therefore defined as
\begin{equation}
	Q(\widehat{\mathbf{P}})
	=
	\frac{
		J(\widehat{\mathbf{P}},\mathbf{U})
	}{
		J_{\max}
	}.
\end{equation}
Hence,
\begin{equation}
	0\leq Q(\widehat{\mathbf{P}})\leq1.
\end{equation}

For the uniform distribution,
\begin{equation}
	Q(\widehat{\mathbf{P}})=0,
\end{equation}
indicating complete equilibrium or maximal disorder. Conversely, $Q(\widehat{\mathbf{P}})=1$ when the entire probability mass is concentrated on a single ordinal pattern, corresponding to the maximum departure from the uniform distribution.

The same Jensen--Shannon-based disequilibrium is used throughout this study to construct the Shannon, Tsallis, R\'enyi, and Fisher-based statistical complexity measures.


\subsubsection{Statistical Complexity Measures}
\label{complexity-measures}

Entropy measures quantify global uncertainty, whereas Fisher information quantifies local changes in the ordinal-pattern distribution. Neither type of measure alone explicitly incorporates the departure of the observed distribution from the uniform distribution. Statistical complexity measures combine these complementary characteristics.

Following the framework of Lamberti et al.~\cite{lamberti2004intensive}, the statistical complexity is constructed by combining an information measure with the normalized disequilibrium $Q(\widehat{\mathbf{P}})$. This construction produces low complexity for distributions approaching complete disorder and also distinguishes highly concentrated distributions from distributions exhibiting intermediate levels of organization.

The Shannon statistical complexity is defined as
\begin{equation}
	C_{S}(\widehat{\mathbf{P}})
	=
	H_{S}(\widehat{\mathbf{P}})
	Q(\widehat{\mathbf{P}}).
\end{equation}

Similarly, the Tsallis statistical complexity is
\begin{equation}
	C_{T}^{(q)}(\widehat{\mathbf{P}})
	=
	H_{T}^{(q)}(\widehat{\mathbf{P}})
	Q(\widehat{\mathbf{P}}),
\end{equation}
where $H_{T}^{(q)}(\widehat{\mathbf{P}})$ is the normalized Tsallis entropy of order $q$.

The R\'enyi statistical complexity is defined as
\begin{equation}
	C_{R}^{(q)}(\widehat{\mathbf{P}})
	=
	H_{R}^{(q)}(\widehat{\mathbf{P}})
	Q(\widehat{\mathbf{P}}),
\end{equation}
where $H_{R}^{(q)}(\widehat{\mathbf{P}})$ denotes the normalized R\'enyi entropy of order $q$.

Finally, the Fisher-based complexity is defined by combining the Fisher Information Measure with the same normalized disequilibrium:
\begin{equation}
	C_{F}(\widehat{\mathbf{P}})
	=
	F(\widehat{\mathbf{P}})
	Q(\widehat{\mathbf{P}}).
\end{equation}

Together, the Shannon, Tsallis, and R\'enyi entropy-complexity measures characterize global properties of the ordinal-pattern distribution, whereas the Fisher Information Measure and Fisher-based complexity provide complementary sensitivity to local structure. These quantities are subsequently used as features to characterize temporal changes in the RSAM time series.

\subsection{Thirteen Features from ordinal patterns}
Table~\ref{tab:features} lists the thirteen features specified for this analysis. Here $H_{\mathrm{Fisher}}$ retains the analysis-file notation but denotes Fisher information, not an entropy.

\begin{table}[H]
	\centering
	\caption{Candidate features, grouped by type.}\label{tab:features}
	\begin{tabularx}{\linewidth}{lYr}
		\toprule
		Group & Features & Count\\
		\midrule
		Entropy/information & $H_{\mathrm{Shannon}}$, $H_{\mathrm{R\acute{e}nyi}}$, $H_{\mathrm{Tsallis}}$, $H_{\mathrm{Fisher}}$ & 4\\
		Complexity & $C_{\mathrm{Shannon}}$, $C_{\mathrm{R\acute{e}nyi}}$, $C_{\mathrm{Tsallis}}$, $C_{\mathrm{Fisher}}$ & 4\\
		Disequilibrium & Disequilibrium & 1\\
		Variance & $\mathrm{Var}(H_{\mathrm{Shannon}})$, $\mathrm{Var}(H_{\mathrm{R\acute{e}nyi}})$, $\mathrm{Var}(H_{\mathrm{Tsallis}})$, $\mathrm{Var}(C_{\mathrm{Shannon}})$ & 4\\
		\bottomrule
	\end{tabularx}
\end{table}

The four complexity measures combine their corresponding normalized information measure with the Jensen--Shannon-based normalized disequilibrium defined in Section~\ref{disequilibrium}.

\subsection{Feature Set Selection}\label{sec:feature-selection}

The thirteen features in Table~\ref{tab:features} include four normalized global entropies ($H_{\mathrm{Shannon}}$, $H_{\mathrm{R\acute{e}nyi}}$, $H_{\mathrm{Tsallis}}$, and the Fisher Information Measure, here labeled $H_{\mathrm{Fisher}}$ for table consistency although it is not an entropy), four corresponding statistical complexities, one disequilibrium value, and four measures of within-window variance. Several of these descriptors are constructed from the same underlying ordinal-pattern distribution and are therefore expected to covary: the three classical normalized entropies ($H_{\mathrm{Shannon}}$, $H_{\mathrm{R\acute{e}nyi}}$, $H_{\mathrm{Tsallis}}$) all summarize the same global disorder of $\widehat{\mathbf{P}}$ under different weighting schemes, and their corresponding complexities share the common disequilibrium factor $Q(\widehat{\mathbf{P}})$ (Section~\ref{complexity-measures}). Figures~\ref{fig:corr_wiz} and~\ref{fig:corr_wsrz} (Section~\ref{sec:corr-results}) confirm this expected redundancy directly, at both stations.

To evaluate whether this redundancy can be reduced without discarding useful information, three nested feature sets are compared throughout the clustering and classification analysis:

\begin{itemize}
	\item \textbf{Thirteen features}: the full descriptor set in Table~\ref{tab:features}, retained as the reference case against which reduced sets are compared.
	\item \textbf{Seven features}: the Fisher Information Measure, Fisher-based complexity, disequilibrium, and the four variance descriptors ($\mathrm{Var}(H_{\mathrm{Shannon}})$, $\mathrm{Var}(C_{\mathrm{Shannon}})$, $\mathrm{Var}(H_{\mathrm{R\acute{e}nyi}})$, $\mathrm{Var}(H_{\mathrm{Tsallis}})$). This set drops the three classical normalized-entropy/complexity pairs (Shannon, R\'enyi, Tsallis), which Figures~\ref{fig:corr_wiz} and~\ref{fig:corr_wsrz} show to be the most strongly intercorrelated group, while retaining a local-structure descriptor (Fisher), a global departure-from-uniformity descriptor (disequilibrium), and four measures of short-term variability that are not directly implied by the level-based descriptors.
	\item \textbf{Four features}: the variance descriptors alone, testing whether the within-window variability of the entropy trajectory carries classification-relevant information on its own, independent of the corresponding level-based (mean-entropy) descriptors.
\end{itemize}

The seven- and four-feature sets are compared for K-means, DBSCAN, and SVM; the thirteen-feature set is additionally included for SVM, where it is computationally practical to test the full descriptor set directly. This comparison is descriptive rather than a formal feature-selection procedure: it reports how clustering quality and classification performance change with feature-set size, rather than selecting a single ``optimal'' subset for the final analysis.

\subsection{VAL Reference Classes}\label{sec:val}

Volcanic Alert Level (VAL) records were obtained from the GeoNet data repository~\cite{GeoNetVAL}. These records describe assessed volcanic status and are used as an external activity reference, rather than as direct annotations of individual seismic events. Each record includes the period start and end, the numeric VAL, the VAL system version, and the official activity description. Period boundaries are provided in both UTC and New Zealand local time (NZST or NZDT).

The analysis uses VAL System Version 3 records overlapping the study period. Table~\ref{tab:val_v3} presents the official six-level classification.

\begin{table}[H]
	\centering
	\caption{Official GeoNet Volcanic Alert Level System
		Version 3 definitions~\cite{GeoNetVAL}.}
	\label{tab:val_v3}
	\begin{tabularx}{\linewidth}{lY}
		\toprule
		VAL & Volcanic activity \\
		\midrule
		0 & No volcanic unrest \\
		1 & Minor volcanic unrest \\
		2 & Moderate to heightened volcanic unrest \\
		3 & Minor volcanic eruption \\
		4 & Moderate volcanic eruption \\
		5 & Major volcanic eruption \\
		\bottomrule
	\end{tabularx}
\end{table}	

For clustering comparisons and supervised classification, the relevant Version 3 activity descriptions are grouped into three classes (Table~\ref{tab:val_mapping}). Class 1 	corresponds to minor volcanic unrest and is termed ``Background unrest'' in the analysis code. This label does not imply an absence of volcanic unrest. Class 2 represents moderate to heightened volcanic unrest, while Class 3 combines minor and moderate volcanic eruption. This mapping does not assign VAL 0 and VAL 5 a class.

\begin{table}[H]
	\centering
	\caption{Three-class grouping used in the analysis.
		The class numbers are study-specific labels, distinct
		from the official VAL scale.}
	\label{tab:val_mapping}
	\begin{tabularx}{\linewidth}{llY}
		\toprule
		Class & VAL & Description \\
		\midrule
		1 & 1 & Minor volcanic unrest
		(background unrest in the analysis) \\
		2 & 2 & Moderate to heightened volcanic unrest \\
		3 & 3--4 & Minor or moderate volcanic eruption \\
		\bottomrule
	\end{tabularx}
\end{table}	

For each weekly window identifies overlapping VAL periods and assigns the highest mapped class. An overlap occurs when the VAL period starts strictly before the window ends and ends strictly after the window starts. Windows with non-finite values in the selected features or without a mapped activity label are excluded. A window assigned to Class 3 therefore overlaps an eruption-labeled VAL period.

Figures~\ref{fig:val_system_sankey} and~\ref{fig:val_activity_sankey} illustrate the relationships between VAL levels, their official descriptions, and the three analysis classes. Flow widths represent the number of recorded VAL periods in the filtered dataset, rather than their durations or the number of weekly windows.

\begin{figure}[H]
	\centering
	\includegraphics[
	width=\linewidth,
	height=0.71\textheight,	keepaspectratio]
	{\figpath/VAL_System_Description_Sankey.pdf}
	\caption{Relationships between VAL levels and official
		activity descriptions in the filtered Version 3 record.
		Flow widths represent the number of recorded VAL periods.}
	\label{fig:val_system_sankey}
\end{figure}

\begin{figure}[H]
	\centering
	\includegraphics[
	width=\linewidth,
	height=0.71\textheight,
	keepaspectratio]{\figpath/VAL_Activity_Class_Sankey.pdf}
	\caption{Mapping from VAL levels to the three analysis
		classes. Flow widths represent numbers of recorded VAL
		periods, rather than their durations.}
	\label{fig:val_activity_sankey}
\end{figure}


\subsection{K-means Clustering}\label{sec:methods-kmeans}

K-means partitioning is applied to the standardized feature matrix for each station and feature set, using Euclidean distance, $k=3$ clusters, and 50 random starts (a fixed seed is used throughout, so results are reproducible). The purpose is to test whether an unsupervised three-group partition of the ordinal-pattern feature space resembles the three VAL reference classes, not to force agreement with them: cluster labels are assigned without reference to VAL, and are only compared with the VAL classes afterward.

Cluster quality is reported two ways. Internally, the between-cluster to total sum-of-squares ratio and the mean silhouette width summarize how compact and well separated the three clusters are, independent of any external reference. Externally, the Adjusted Rand Index (ARI) against the VAL-derived reference labels measures agreement with the VAL classes while correcting for chance: $\mathrm{ARI}=1$ indicates perfect agreement, $\mathrm{ARI}=0$ indicates agreement no better than a random partition of the same sizes, and negative values indicate worse than random agreement.

\subsection{DBSCAN Clustering}\label{sec:methods-dbscan}

DBSCAN groups nearby feature vectors into density-connected clusters and leaves points that do not satisfy its density conditions unassigned (``noise''). Rather than fixing $\varepsilon$ and $\mathrm{minPts}$ in advance, a grid search is performed separately for each station and feature set over
\begin{align*}
	\varepsilon &\in \{0.05, 0.10, \ldots, 5.00\}, \\
	\mathrm{minPts} &\in \{3, 4, \ldots, 20\},
\end{align*}
and every combination producing exactly three non-noise clusters -- a number chosen to match the three VAL reference classes, rather than the conventional DBSCAN use of one dense cluster plus noise -- is retained as a candidate. Among these candidates, the combination with the fewest noise points is selected, with ties broken first by proximity of $\varepsilon$ to a reference value of 1.1 and then by the smaller $\mathrm{minPts}$. Under the selected parameters, points that remain outside all three clusters are still reported as noise and are treated as candidate anomalous windows; they are not automatically eruption precursors, and membership of a dense cluster does not establish a physically normal state. As with K-means, the resulting three-cluster partition is compared with the VAL classes only after fitting.

\subsection{Clustering Quality Across Feature Sets}\label{sec:methods-clustering-comparison}

Because the DBSCAN parameter search is comparatively expensive to repeat, the comparison of clustering quality across feature sets is restricted to the seven- and four-feature sets (Section~\ref{sec:feature-selection}) for both K-means and DBSCAN, at each station. For each combination, the average silhouette width and the Adjusted Rand Index against the VAL classes are recorded, giving four station-level comparisons (K-means and DBSCAN, seven and four features).

\subsection{Support Vector Machine Classification}\label{sec:methods-svm}

A radial-basis-function (RBF) support vector machine (SVM) addresses the corresponding supervised, three-class classification problem (background unrest, moderate-to-heightened unrest, eruption). Unlike the clustering methods above, the SVM is evaluated by cross-validation rather than by comparison with an unsupervised partition.

For each of the three feature sets (thirteen, seven, and four descriptors; Section~\ref{sec:feature-selection}) and four class-imbalance treatments -- no balancing, random oversampling of the minority classes, random undersampling of the majority class, and SMOTE (synthetic minority oversampling, using each minority-class observation's $k=5$ nearest same-class neighbours) -- a $k$-fold cross-validation is performed on the standardized feature matrix, with $k=\min(5,\,n_{\min})$ and $n_{\min}$ the size of the smallest class. Because the eruption class contains only four weeks at each station, this reduces $k$ to a four-fold split throughout. Class balancing is applied separately within each training fold only, never to the held-out fold, so held-out predictions are never evaluated on synthetic or duplicated observations.

Within each training fold, after balancing, class weights inversely proportional to class frequency are computed as
\begin{equation}
	w_k = \frac{N_{\mathrm{train}}}{3\,n_k},
\end{equation}
where $N_{\mathrm{train}}$ is the (possibly balanced) number of training windows and $n_k$ is the number of training windows in class $k$, so that all three classes contribute comparably to the fitted decision boundary regardless of their post-balancing sample sizes. An RBF kernel is used with a fixed cost $C=1$ and kernel parameter $\gamma = 1/p$, where $p$ is the number of features in the set being evaluated; cost and $\gamma$ are held constant across all comparisons rather than tuned by grid search, so that the reported differences reflect the feature set and balancing method rather than SVM hyperparameter selection.

For each station, feature set, and balancing method, the resulting out-of-fold predictions give an overall accuracy and, following the per-class recall, precision, and F1 definitions in Section~\ref{sec:methods-metrics} below, a one-versus-rest balanced accuracy (the mean per-class recall). A separate baseline run -- seven features, no balancing -- is used for the main confusion matrix and PCA classification figure, so that a single reference configuration is available for direct comparison with the K-means and DBSCAN results in Section~\ref{sec:results}.

\subsubsection{Per-Class Metrics}\label{sec:methods-metrics}

For a given set of predictions, per-class recall, precision, and F1 are calculated in the usual way for each class $k$: recall is the fraction of true class-$k$ windows correctly predicted, precision is the fraction of predicted class-$k$ windows that are correctly labeled, and F1 is their harmonic mean (defined as zero when a class receives no correct predictions and either recall or precision is undefined). Given only three or four eruption-labeled windows per station, raw correct/total counts are reported alongside these ratios throughout, since ``2 of 4'' is more informative than ``50\%'' at this sample size.

Since hyperparameter values are fixed rather than tuned, and the same weekly windows are used for both fitting and out-of-fold evaluation across all folds, the reported cross-validated performance should not be interpreted as assessment on an independent external test set. In addition, because ordinal-pattern features are computed from overlapping raw observations at adjacent weekly steps, windows assigned to different folds may not be fully independent. A blocked or event-separated validation strategy would provide a more stringent assessment of generalization under this temporal dependence and is identified as an important direction for future validation, together with a hyperparameter-tuned SVM configuration.

\section{Results}\label{sec:results}
\subsection{Data Coverage and VAL Class Counts}
After correcting a VAL date-parsing issue that had silently dropped a four-month span of VAL records at both stations (Section~\ref{sec:val}), the modeling output contains 335 WIZ and 333 WSRZ windows. Table~\ref{tab:label_counts} shows the corrected counts. The code calls Class 1 ``Background unrest''; this corresponds to the VAL description ``Minor volcanic unrest'', not the absence of unrest. Each station now contributes four eruption-labeled (Class 3) windows.

\begin{table}[H]
	\centering\small
	\caption{Corrected weekly class counts.}\label{tab:label_counts}
	\begin{tabularx}{\linewidth}{Yrrrr}
		\toprule
		Station & Class 1 & Class 2 & Class 3 & Total\\
		\midrule
		WIZ & 320 & 11 & 4 & 335\\
		WSRZ & 318 & 11 & 4 & 333\\
		\bottomrule
	\end{tabularx}
\end{table}

\subsection{Entropy-Complexity Planes}
Figures~\ref{fig:hc_wiz} and~\ref{fig:hc_wsrz} retain the four supplied panels for each station. Shannon, Tsallis, and R\'enyi panels show concentrated curved relationships, while the Fisher panels show a different increasing association. Several labeled windows lie towards the edges of these distributions. WSRZ weeks 146--151 recur among the labeled points, and week 151 is labeled in several WIZ panels. Match week indices to dates before attributing event associations; they do not identify eruption windows. Colors encode week index rather than VAL class. The Fisher panel's horizontal coordinate denotes Fisher information despite the shared $H$ axis notation.

\begin{figure}[H]
	\centering
	\includegraphics[width=\linewidth,height=0.71\textheight,keepaspectratio]{\figpath/HC_Plane_WIZ_D5.pdf}
	\caption{WIZ ordinal-feature planes for $D=5$, with Shannon, Tsallis, R\'enyi, and Fisher panels. Colors and selected week labels are retained from the supplied figure.}\label{fig:hc_wiz}
\end{figure}

\begin{figure}[H]
	\centering
	\includegraphics[width=\linewidth,height=0.71\textheight,keepaspectratio]{\figpath/HC_Plane_WSRZ_D5.pdf}
	\caption{WSRZ ordinal-feature planes for $D=5$. The panels retain their original scales, weak colors, labels, and guide lines.}\label{fig:hc_wsrz}
\end{figure}

Table~\ref{tab:hc_extremes_both_stations} compares the selected extreme windows for WIZ and WSRZ. Week 151, covering 17--24 December 2016, is selected in the Shannon, R\'enyi, and Tsallis planes at both stations. Among the
listed windows, it has the lowest entropy and highest complexity for these three measures at each station. The selected Fisher windows differ between stations: weeks 15, 68, 255, and 330 at WIZ, and weeks 116, 126, 147, and 151 at WSRZ. These selections describe extremes in the feature planes. 

\begin{table}[H]
	\centering
	\small
	\setlength{\tabcolsep}{4pt}
	\caption{Selected extreme windows for WIZ and WSRZ in the
		ordinal-feature planes ($D=5$).}
	\label{tab:hc_extremes_both_stations}
	\begin{tabular}{llrllrr}
		\toprule
		Measure & Station & Week & Start date & End date & $H$ & $C$ \\
		\midrule
		Shannon & WIZ & 91 & 2015-12-23 & 2015-12-30 & 0.91698 & 0.12138 \\
		Shannon & WIZ & 99 & 2016-02-09 & 2016-02-16 & 0.91937 & 0.12436 \\
		Shannon & WIZ & 151 & 2016-12-17 & 2016-12-24 & 0.89072 & 0.13647 \\
		Shannon & WIZ & 255 & 2018-09-02 & 2018-09-09 & 0.92004 & 0.12722 \\
		\addlinespace[2pt]
		Shannon & WSRZ & 146 & 2016-11-17 & 2016-11-24 & 0.88324 & 0.16222 \\
		Shannon & WSRZ & 147 & 2016-11-23 & 2016-11-30 & 0.87397 & 0.17108 \\
		Shannon & WSRZ & 150 & 2016-12-11 & 2016-12-18 & 0.88377 & 0.15582 \\
		Shannon & WSRZ & 151 & 2016-12-17 & 2016-12-24 & 0.85154 & 0.18292 \\
		\midrule
		Tsallis & WIZ & 91 & 2015-12-23 & 2015-12-30 & 0.96421 & 0.12319 \\
		Tsallis & WIZ & 99 & 2016-02-09 & 2016-02-16 & 0.96734 & 0.12629 \\
		Tsallis & WIZ & 104 & 2016-03-10 & 2016-03-17 & 0.96637 & 0.12472 \\
		Tsallis & WIZ & 151 & 2016-12-17 & 2016-12-24 & 0.94477 & 0.13971 \\
		\addlinespace[2pt]
		Tsallis & WSRZ & 146 & 2016-11-17 & 2016-11-24 & 0.94808 & 0.16807 \\
		Tsallis & WSRZ & 147 & 2016-11-23 & 2016-11-30 & 0.94233 & 0.17804 \\
		Tsallis & WSRZ & 150 & 2016-12-11 & 2016-12-18 & 0.94699 & 0.16115 \\
		Tsallis & WSRZ & 151 & 2016-12-17 & 2016-12-24 & 0.92604 & 0.19200 \\
		\midrule
		R\'enyi & WIZ & 91 & 2015-12-23 & 2015-12-30 & 0.87270 & 0.11150 \\
		R\'enyi & WIZ & 104 & 2016-03-10 & 2016-03-17 & 0.87938 & 0.11349 \\
		R\'enyi & WIZ & 151 & 2016-12-17 & 2016-12-24 & 0.81697 & 0.12081 \\
		R\'enyi & WIZ & 255 & 2018-09-02 & 2018-09-09 & 0.88691 & 0.11837 \\
		\addlinespace[2pt]
		R\'enyi & WSRZ & 146 & 2016-11-17 & 2016-11-24 & 0.82596 & 0.14642 \\
		R\'enyi & WSRZ & 147 & 2016-11-23 & 2016-11-30 & 0.81047 & 0.15313 \\
		R\'enyi & WSRZ & 150 & 2016-12-11 & 2016-12-18 & 0.82296 & 0.14005 \\
		R\'enyi & WSRZ & 151 & 2016-12-17 & 2016-12-24 & 0.76952 & 0.15955 \\
		\midrule
		Fisher & WIZ & 15 & 2014-09-23 & 2014-09-30 & 0.10010 & 0.01311 \\
		Fisher & WIZ & 68 & 2015-08-07 & 2015-08-14 & 0.09764 & 0.01285 \\
		Fisher & WIZ & 255 & 2018-09-02 & 2018-09-09 & 0.09643 & 0.01287 \\
		Fisher & WIZ & 330 & 2019-11-26 & 2019-12-03 & 0.10240 & 0.01298 \\
		\addlinespace[2pt]
		Fisher & WSRZ & 116 & 2016-05-21 & 2016-05-28 & 0.11279 & 0.01403 \\
		Fisher & WSRZ & 126 & 2016-07-20 & 2016-07-27 & 0.11506 & 0.01568 \\
		Fisher & WSRZ & 147 & 2016-11-23 & 2016-11-30 & 0.12750 & 0.02409 \\
		Fisher & WSRZ & 151 & 2016-12-17 & 2016-12-24 & 0.11120 & 0.02306 \\
		\bottomrule
	\end{tabular}
\end{table}


\subsection{Feature Correlations}\label{sec:corr-results}
Figures~\ref{fig:corr_wiz} and~\ref{fig:corr_wsrz} show Spearman rank correlations among all thirteen candidate features (Table~\ref{tab:features}) at each station, ordered by hierarchical clustering of the correlation matrix. The three classical normalized entropy/complexity pairs -- Shannon, R\'enyi, and Tsallis -- form a tightly correlated block at both stations, with $|\rho|$ close to 1 between $H_{\mathrm{R\acute{e}nyi}}$ and $H_{\mathrm{Tsallis}}$, and between the corresponding complexity measures; this is expected, since all three summarize the same ordinal-pattern distribution under different weighting schemes (Sections~\ref{disequilibrium} and~\ref{complexity-measures}). The Fisher Information Measure, Fisher-based complexity, and disequilibrium are likewise strongly correlated with one another ($\rho\gtrsim0.9$), and the four variance descriptors form their own tightly correlated block, separate from the level-based (mean) descriptors. These patterns motivate the seven- and four-feature sets described in Section~\ref{sec:feature-selection}: the seven-feature set removes the most redundant block (the three classical entropy/complexity pairs) while retaining one representative of the Fisher-related block, the disequilibrium measure, and the four variance descriptors; the four-feature set retains only the variance block. The correlation structure is qualitatively similar between WIZ and WSRZ, indicating that this redundancy is a property of the feature construction rather than of either station's record specifically.

\begin{figure}[H]
	\centering
	\includegraphics[width=\linewidth,height=0.71\textheight,keepaspectratio]
	{\figpath/WIZ_13_Feature_Correlogram.pdf}
	\caption{Spearman rank correlations among all thirteen WIZ features.}
	\label{fig:corr_wiz}
\end{figure}

\begin{figure}[H]
	\centering
	\includegraphics[width=\linewidth,height=0.71\textheight,keepaspectratio]
	{\figpath/WSRZ_13_Feature_Correlogram.pdf}
	\caption{Spearman rank correlations among all thirteen WSRZ features.}
	\label{fig:corr_wsrz}
\end{figure}

\subsection{Clustering Results}\label{sec:clustering-results}

\subsubsection{K-means}

Table~\ref{tab:kmeans} cross-tabulates the K-means partition (seven features) against the VAL reference classes. At both stations, the three K-means clusters split the dominant background class roughly in two while distributing the small moderate- and eruption-class counts across the same two clusters, rather than isolating either as a distinct group; a third, much smaller cluster picks up a handful of windows without a clear correspondence to VAL class. The between-cluster to total sum-of-squares ratio is 0.560 (WIZ) and 0.520 (WSRZ), and the mean silhouette width is 0.454 (WIZ) and 0.357 (WSRZ), indicating a moderately well-separated three-group structure in feature space -- but, as Table~\ref{tab:kmeans} shows, one that does not align with the VAL classes.

\begin{table}[H]
	\centering\small
	\caption{K-means partition (seven features) versus VAL class.}\label{tab:kmeans}
	\begin{tabularx}{\linewidth}{Ylrrrr}
		\toprule
		Station & VAL class & Cluster 1 & Cluster 2 & Cluster 3 & Total\\
		\midrule
		WIZ & 1 (Background) & 111 & 207 & 2 & 320\\
		WIZ & 2 (Moderate) & 4 & 7 & 0 & 11\\
		WIZ & 3 (Eruption) & 3 & 1 & 0 & 4\\
		\addlinespace[2pt]
		WSRZ & 1 (Background) & 158 & 156 & 4 & 318\\
		WSRZ & 2 (Moderate) & 5 & 6 & 0 & 11\\
		WSRZ & 3 (Eruption) & 2 & 2 & 0 & 4\\
		\bottomrule
	\end{tabularx}
\end{table}

\begin{figure}[H]
	\centering
	\includegraphics[width=\linewidth,height=0.6\textheight,keepaspectratio]{\figpath/WIZ_Kmeans_Cluster_Plot.pdf}
	\caption{WIZ: K-means clusters (seven features) in PCA coordinates (PC1 63.8\%, PC2 24.3\% of standardized-feature variance).}\label{fig:kmeans_WIZ}
\end{figure}

\begin{figure}[H]
	\centering
	\includegraphics[width=\linewidth,height=0.6\textheight,keepaspectratio]{\figpath/WSRZ_Kmeans_Cluster_Plot.pdf}
	\caption{WSRZ: K-means clusters (seven features) in PCA coordinates (PC1 65.3\%, PC2 24.8\% of standardized-feature variance).}\label{fig:kmeans_WSRZ}
\end{figure}

\subsubsection{DBSCAN}

The grid search (Section~\ref{sec:methods-dbscan}) selects $\varepsilon=1$, $\mathrm{minPts}=4$ for WIZ and $\varepsilon=1.05$, $\mathrm{minPts}=3$ for WSRZ. Under these parameters, 5.07\% of WIZ windows (17 of 335) and 4.50\% of WSRZ windows (15 of 333) are left as noise; the remainder are split across three clusters that are, as with K-means, dominated by a single large cluster (Table~\ref{tab:dbscan}). Mean silhouette width is 0.325 (WIZ) and 0.207 (WSRZ) -- lower than for K-means at both stations, consistent with DBSCAN's density criterion producing less balanced cluster sizes than a forced three-way partition. No eruption-labeled window is left as noise at either station: all four WIZ eruption windows and all four WSRZ eruption windows fall inside a cluster, split between Cluster 1 and Cluster 2.

\begin{table}[H]
	\centering\small
	\caption{DBSCAN partition (seven features) versus VAL class.}\label{tab:dbscan}
	\begin{tabularx}{\linewidth}{Ylrrrrr}
		\toprule
		Station & VAL class & Noise & Cluster 1 & Cluster 2 & Cluster 3 & Total\\
		\midrule
		WIZ & 1 (Background) & 16 & 294 & 2 & 8 & 320\\
		WIZ & 2 (Moderate) & 1 & 10 & 0 & 0 & 11\\
		WIZ & 3 (Eruption) & 0 & 3 & 1 & 0 & 4\\
		\addlinespace[2pt]
		WSRZ & 1 (Background) & 14 & 297 & 3 & 4 & 318\\
		WSRZ & 2 (Moderate) & 1 & 9 & 0 & 1 & 11\\
		WSRZ & 3 (Eruption) & 0 & 4 & 0 & 0 & 4\\
		\bottomrule
	\end{tabularx}
\end{table}

\begin{figure}[H]
	\centering
	\includegraphics[width=\linewidth,height=0.6\textheight,keepaspectratio]{\figpath/WIZ_DBSCAN_Cluster_Plot.pdf}
	\caption{WIZ: DBSCAN clusters and noise (seven features) in PCA coordinates.}\label{fig:dbscan_WIZ}
\end{figure}

\begin{figure}[H]
	\centering
	\includegraphics[width=\linewidth,height=0.6\textheight,keepaspectratio]{\figpath/WSRZ_DBSCAN_Cluster_Plot.pdf}
	\caption{WSRZ: DBSCAN clusters and noise (seven features) in PCA coordinates.}\label{fig:dbscan_WSRZ}
\end{figure}

\subsubsection{Clustering Quality Across Feature Sets}

Table~\ref{tab:cluster_quality} compares the seven- and four-feature sets for both methods, using average silhouette width and Adjusted Rand Index (ARI) against the VAL classes. ARI is close to zero, and in two of four WSRZ comparisons slightly negative, for every combination: no configuration recovers a partition that agrees with the VAL classes better than chance to a practically meaningful degree. The four-feature (variance-only) set gives a higher silhouette width than the seven-feature set for K-means at both stations, but a lower silhouette width than the seven-feature set for DBSCAN, particularly at WSRZ, where the four-feature DBSCAN silhouette is effectively zero. Reducing to four features therefore changes the internal geometry of the clusters without improving external agreement with VAL class in either direction.

\begin{table}[H]
	\centering\small
	\caption{Clustering quality by feature set. ARI: Adjusted Rand Index against VAL class.}\label{tab:cluster_quality}
	\begin{tabularx}{\linewidth}{Ylrrr}
		\toprule
		Station & Feature set & Method & Avg. silhouette & ARI\\
		\midrule
		WIZ & 7 features & K-means & 0.454 & 0.010\\
		WIZ & 7 features & DBSCAN & 0.325 & 0.035\\
		WIZ & 4 features & K-means & 0.471 & 0.018\\
		WIZ & 4 features & DBSCAN & 0.114 & 0.052\\
		\addlinespace[2pt]
		WSRZ & 7 features & K-means & 0.357 & $-0.002$\\
		WSRZ & 7 features & DBSCAN & 0.207 & 0.049\\
		WSRZ & 4 features & K-means & 0.467 & $-0.002$\\
		WSRZ & 4 features & DBSCAN & 0.001 & 0.026\\
		\bottomrule
	\end{tabularx}
\end{table}

\begin{figure}[H]
	\centering
	\includegraphics[width=\linewidth,height=0.55\textheight,keepaspectratio]{\figpath/WIZ_Clustering_FeatureSet_Comparison.pdf}
	\caption{WIZ: silhouette width and ARI by clustering method and feature set.}\label{fig:cluster_quality_WIZ}
\end{figure}

\begin{figure}[H]
	\centering
	\includegraphics[width=\linewidth,height=0.55\textheight,keepaspectratio]{\figpath/WSRZ_Clustering_FeatureSet_Comparison.pdf}
	\caption{WSRZ: silhouette width and ARI by clustering method and feature set.}\label{fig:cluster_quality_WSRZ}
\end{figure}

\subsection{SVM Classification Results}\label{sec:svm-results}

\subsubsection{Baseline (Seven Features, No Balancing)}

Table~\ref{tab:svm_confusion} gives the out-of-fold confusion matrix for the baseline configuration (seven features, no balancing, three-fold cross-validation). Overall accuracy is 0.627 at both stations, but this is driven almost entirely by the background class: one-versus-rest balanced accuracy is only 0.365 (WIZ) and 0.530 (WSRZ). WIZ correctly recovers none of its four eruption windows; WSRZ correctly recovers two of four.

\begin{table}[H]
	\centering\small
	\caption{Baseline SVM (seven features, no balancing) out-of-fold confusion matrix.}\label{tab:svm_confusion}
	\begin{tabularx}{\linewidth}{Ylrrr}
		\toprule
		Station & True class & Pred. 1 & Pred. 2 & Pred. 3\\
		\midrule
		WIZ & 1 (Background, $n=320$) & 205 & 87 & 28\\
		WIZ & 2 (Moderate, $n=11$) & 4 & 5 & 2\\
		WIZ & 3 (Eruption, $n=4$) & 3 & 1 & 0\\
		\addlinespace[2pt]
		WSRZ & 1 (Background, $n=318$) & 202 & 96 & 20\\
		WSRZ & 2 (Moderate, $n=11$) & 6 & 5 & 0\\
		WSRZ & 3 (Eruption, $n=4$) & 2 & 0 & 2\\
		\bottomrule
	\end{tabularx}
\end{table}

At WIZ, 28 background windows and 2 moderate windows are predicted as Class 3, alongside 0 of 4 true eruption windows: the eruption-class predictions in this configuration are entirely false positives. At WSRZ, both correctly predicted eruption windows occur alongside no misclassified moderate windows into Class 3, giving an eruption recall of $2/4=0.50$ and a precision of $2/2=1.00$ for this baseline pass -- notably higher than at WIZ, though still based on only two correctly recovered windows.

\subsubsection{Feature Set and Balancing Method Comparison}

Table~\ref{tab:svm_comparison} reports accuracy, one-versus-rest balanced accuracy, and eruption-class recall across all three feature sets and four balancing methods, at both stations; Figures~\ref{fig:svm_WIZ} and~\ref{fig:svm_WSRZ} display the same comparison graphically, together with eruption recall.

At WSRZ, eruption recall reaches $2/4=0.50$ under nine of the twelve feature-set/balancing combinations, including the unbalanced thirteen- and seven-feature configurations; only the four-feature set with no balancing or with undersampling falls to $1/4=0.25$. Balanced accuracy at WSRZ is highest for the thirteen-feature SMOTE configuration (0.560) and lowest for the four-feature undersampled configuration (0.341).

At WIZ, by contrast, eruption recall is $0/4$ under eleven of the twelve combinations; the sole exception is the seven-feature, undersampled configuration, which recovers $1/4=0.25$ of eruption windows but only at a substantial cost elsewhere, with overall accuracy falling to 0.254 (from 0.648 unbalanced) and balanced accuracy falling to 0.200 (below the unbalanced baseline of 0.343). No feature set or balancing method tested recovers more than one of the four WIZ eruption windows.

\begin{table}[H]
	\centering\small
	\setlength{\tabcolsep}{4pt}
	\caption{SVM accuracy, one-versus-rest balanced accuracy, and eruption recall (correct of 4) by feature set and balancing method.}\label{tab:svm_comparison}
	\begin{tabularx}{\linewidth}{Ylllrrr}
		\toprule
		Station & Feature set & Balancing & Accuracy & Bal. accuracy & Eruption recall\\
		\midrule
		WIZ & 13 features & No balancing & 0.651 & 0.432 & 0/4\\
		WIZ & 13 features & Oversample & 0.618 & 0.391 & 0/4\\
		WIZ & 13 features & Undersample & 0.373 & 0.247 & 0/4\\
		WIZ & 13 features & SMOTE & 0.687 & 0.386 & 0/4\\
		WIZ & 7 features & No balancing & 0.648 & 0.343 & 0/4\\
		WIZ & 7 features & Oversample & 0.588 & 0.293 & 0/4\\
		WIZ & 7 features & Undersample & 0.254 & 0.200 & 1/4\\
		WIZ & 7 features & SMOTE & 0.666 & 0.320 & 0/4\\
		WIZ & 4 features & No balancing & 0.412 & 0.290 & 0/4\\
		WIZ & 4 features & Oversample & 0.504 & 0.322 & 0/4\\
		WIZ & 4 features & Undersample & 0.466 & 0.338 & 0/4\\
		WIZ & 4 features & SMOTE & 0.510 & 0.354 & 0/4\\
		\addlinespace[2pt]
		WSRZ & 13 features & No balancing & 0.592 & 0.547 & 2/4\\
		WSRZ & 13 features & Oversample & 0.688 & 0.551 & 2/4\\
		WSRZ & 13 features & Undersample & 0.294 & 0.384 & 2/4\\
		WSRZ & 13 features & SMOTE & 0.715 & 0.560 & 2/4\\
		WSRZ & 7 features & No balancing & 0.637 & 0.533 & 2/4\\
		WSRZ & 7 features & Oversample & 0.673 & 0.516 & 2/4\\
		WSRZ & 7 features & Undersample & 0.369 & 0.440 & 2/4\\
		WSRZ & 7 features & SMOTE & 0.706 & 0.557 & 2/4\\
		WSRZ & 4 features & No balancing & 0.514 & 0.379 & 1/4\\
		WSRZ & 4 features & Oversample & 0.508 & 0.400 & 2/4\\
		WSRZ & 4 features & Undersample & 0.237 & 0.341 & 1/4\\
		WSRZ & 4 features & SMOTE & 0.508 & 0.488 & 2/4\\
		\bottomrule
	\end{tabularx}
\end{table}

\begin{figure}[H]
	\centering
	\includegraphics[width=\linewidth,height=0.71\textheight,keepaspectratio]{\figpath/WIZ_SVM_FeatureSet_Balance_Plot.pdf}
	\caption{WIZ: accuracy, balanced accuracy, and eruption recall by feature set and balancing method.}\label{fig:svm_WIZ}
\end{figure}

\begin{figure}[H]
	\centering
	\includegraphics[width=\linewidth,height=0.71\textheight,keepaspectratio]{\figpath/WSRZ_SVM_FeatureSet_Balance_Plot.pdf}
	\caption{WSRZ: accuracy, balanced accuracy, and eruption recall by feature set and balancing method.}\label{fig:svm_WSRZ}
\end{figure}

\section{Discussion}

The corrected results show that unsupervised structure in the ordinal-pattern feature space and the VAL eruption classes describe different aspects of these data. Both K-means and DBSCAN recover three-group partitions with reasonable internal geometry (silhouette widths up to about 0.47), but none of the eight clustering configurations tested (two methods $\times$ two feature sets $\times$ two stations) reaches an Adjusted Rand Index practically distinguishable from chance agreement with the VAL classes. Unlike the earlier fixed-parameter DBSCAN configuration, which isolated a small noise set entirely outside the eruption class, the present three-cluster DBSCAN search places every eruption-labeled window inside a cluster at both stations; the dominant cluster structure in this feature space is therefore not primarily organized around the eruption/non-eruption distinction, whichever clustering method or feature set is used. This does not mean the ordinal-pattern features are uninformative about VAL class -- the SVM results below show some class-related signal -- only that an unsupervised three-way split of this particular feature space does not reproduce the VAL grouping.

SVM performance is limited and markedly asymmetric between stations. WSRZ recovers two of its four eruption windows under most feature-set and balancing combinations, including with no balancing at all, and does so at WSRZ's baseline configuration without any false-positive eruption predictions. WIZ, in contrast, fails to recover any eruption window under eleven of twelve tested configurations; the one exception (seven features, undersampled) trades a large amount of overall and background-class accuracy for a single additional true positive, and does not represent a favorable trade-off in absolute terms. Because each station contributes only four eruption-labeled windows, these differences rest on very small counts -- a single window's outcome changes a station's recall by 25 percentage points -- and should not be read as evidence that WSRZ's seismic signal is inherently more classifiable than WIZ's without independent corroboration (e.g., differences in station siting, noise characteristics, or which sub-events fall in each station's four eruption weeks).

The thirteen-feature correlogram (Figures~\ref{fig:corr_wiz} and~\ref{fig:corr_wsrz}) confirms substantial redundancy among the classical Shannon/R\'enyi/Tsallis entropy-complexity descriptors, motivating the seven- and four-feature sets. In practice, this reduction changes performance only modestly and inconsistently: the four-feature (variance-only) set generally lowers both SVM balanced accuracy and DBSCAN cluster quality relative to seven or thirteen features, while improving K-means silhouette width at both stations. No single feature set is uniformly best across methods and stations, which argues against treating any one reduced set as a settled choice without further, method-specific validation.

The reported SVM scores are cross-validated on the full available dataset with fixed (untuned) hyperparameters, not evaluated on an independent external test set, and should be read accordingly. Because ordinal-pattern features are computed from overlapping raw observations at adjacent weekly steps, windows assigned to different folds may not be fully independent; a blocked or event-separated validation strategy, together with a properly tuned hyperparameter search, would give a more stringent assessment of generalization. Station comparisons should also account for the small and unequal eruption-class counts and for differences in station-specific signal properties, rather than attributing every performance difference to physical sensitivity.

A forecasting interpretation would require windows ending before an event, labels appropriate to a future prediction horizon, and evaluation based only on information available at the time. The present results support retrospective feature analysis, comparison of clustering and classification configurations, and candidate-anomaly inspection, but not operational forecasting.

\section{Conclusion}

The corrected analysis uses 335 WIZ and 333 WSRZ weekly windows, each with four eruption-labeled windows. Across three feature sets (thirteen, seven, and four ordinal-pattern descriptors), neither K-means nor DBSCAN recovers a three-group partition that agrees with the VAL activity classes beyond chance (Adjusted Rand Index below 0.06 throughout), although both methods produce internally coherent clusters by silhouette width. A class-weighted SVM, evaluated across the same feature sets and four class-imbalance treatments, reaches a baseline (seven-feature, unbalanced) accuracy of 0.627 at both stations but a one-versus-rest balanced accuracy of only 0.365 (WIZ) and 0.530 (WSRZ); WSRZ recovers two of its four eruption windows under most tested configurations, while WIZ recovers at most one, at a substantial cost elsewhere. These results support exploratory characterization of RSAM feature variability and a systematic comparison of clustering and classification configurations, while reliable eruption recognition -- and, a fortiori, forecasting -- remains unestablished with the present feature sets and validation strategy.

\bibliographystyle{plain}
\bibliography{seismic_references}	
\end{document}